\chapter{引言}\label{chap:intro}

\section{研究背景与意义}\label{sec:background}

期权是金融衍生品市场的核心工具之一。根据国际清算银行（BIS）统计，全球场外衍生品市场名义本金规模超过600万亿美元，其中期权类产品占据重要份额。期权定价的准确性直接影响风险对冲、投资组合管理和监管资本计算，是金融工程的核心问题。

自Black、Scholes与Merton于1973年提出BSM模型\cite{black1973pricing,merton1973theory}以来，期权定价理论经历了从常数波动率到局部波动率（CEV模型\cite{cox1976valuation}）、再到随机波动率（Heston模型\cite{heston1993closed}）的演进。这三大模型分别对应不同的市场假设，在实践中各有适用场景：BSM适用于波动率相对平稳的短期定价；CEV能够捕捉股票市场中波动率随股价下跌而上升的杠杆效应；Heston则通过引入随机波动率过程，更准确地描述波动率微笑（volatility smile）和厚尾收益率分布。

然而，传统数值方法（有限差分、有限元）在求解高维PDE时面临"维数灾难"：网格点数随维度指数增长，Heston这样的二维模型已需要大量计算资源。此外，传统工具要求用户手动指定模型类型和参数，对非专业用户门槛较高——选择合适的定价模型本身就需要量化金融背景，普通用户难以判断在当前市场条件下应使用BSM、CEV还是Heston。

物理信息神经网络（Physics-Informed Neural Networks，PINN）\cite{raissi2019physics}提供了一种新的思路：将PDE残差直接嵌入神经网络的损失函数，无需网格剖分，训练完成后可对任意输入即时推断（单次推断约2ms，远低于有限差分的约500ms）。近年来，PINN已被应用于BSM\cite{dhiman2023pinn}、CEV\cite{bae2024localvol}、Heston\cite{hainaut2024heston,guo2024icpinn}等期权定价模型，展现出良好的精度和效率。

与此同时，大语言模型（Large Language Models，LLM）在自然语言理解和结构化信息提取方面取得了突破性进展。将LLM与科学计算工具结合，使用户能够通过自然语言描述问题并获得计算结果，是人机交互的重要方向\cite{liu2025pdeagent}。在期权定价场景中，LLM的核心价值在于：将"选择合适的定价模型"这一需要专业知识的决策自动化，使非专业用户只需描述期权需求，系统自动完成模型选择和参数提取。

\section{现有工作的局限性}\label{sec:motivation}

综合现有PINN期权定价研究，存在以下主要不足：

\textbf{模型孤立，缺乏统一框架。}现有工作（Dhiman等\cite{dhiman2023pinn}、Bae等\cite{bae2024localvol}、Hainaut等\cite{hainaut2024heston}）均针对单一模型分别训练独立的PINN，三个模型之间无法共享网络权重和训练经验。当需要同时支持多种模型时，必须维护多套独立的训练流程和模型文件，工程复杂度高。

\textbf{交互门槛高，缺乏自然语言接口。}现有PINN工具要求用户手动指定模型类型（BSM/CEV/Heston）和所有数值参数，对非专业用户不友好。不同市场条件下应选用何种模型，需要量化金融背景，普通用户难以判断。

\textbf{Heston模型收敛困难。}Heston PDE是二维问题，包含交叉偏导项$\partial^2 V/\partial S\partial v$，独立训练的PINN在Heston上往往收敛困难（本文实验中独立Heston PINN的训练损失停留在$10^3$量级，无法收敛）。

本文针对上述不足，提出将三大模型统一为单一参数化PINN的方法，并结合LLM路由层构建完整的自然语言期权定价系统。

\section{研究目标与贡献}\label{sec:contributions}

本文的研究目标是构建一个基于PINN与LLM的统一期权定价框架，实现以下目标：

\begin{enumerate}
  \item 设计统一PDE算子，通过软掩码机制将BSM、CEV、Heston三大模型的PDE统一为单一参数化形式，由一个网络同时覆盖三种模型；
  \item 采用加法参数化输出，以Black-Scholes解析解为基线，解决深度虚值区域的数值退化问题；
  \item 对Heston模型采用ICPINN硬约束技术，确保终值条件精确满足；
  \item 引入LLM作为智能路由层，支持中英文自然语言输入，自动完成模型选择和参数提取；
  \item 通过与解析解、有限差分参考解的对比，定量评估统一PINN的精度，并与独立PINN基线对比。
\end{enumerate}

本文的主要贡献如下：

\begin{itemize}
  \item \textbf{统一PDE算子设计}：提出软掩码机制$\text{mask}=\tanh(\xi/0.05)^2$，将三大模型的扩散系数统一为$a=(1-\text{mask})\sigma^2 S^{2\beta}+\text{mask}\cdot v S^2$，实现模型间的连续插值，单一网络覆盖三种定价模型。
  \item \textbf{加法参数化输出}：以Black-Scholes解析解$V_\text{BS}$为基线，网络输出$V=V_\text{BS}+K\cdot\text{net}(x)$，解决了乘法参数化在深度虚值区域$V_\text{BS}\to 0$时的退化问题，同时提供良好的初始化（训练开始时$V\approx V_\text{BS}$）。
  \item \textbf{LLM路由集成}：将DeepSeek LLM与统一PINN求解器结合，通过结构化JSON输出实现可靠的参数提取和模型选择，支持中英文自然语言输入，降低期权定价工具的使用门槛。
  \item \textbf{实验验证}：在BSM（MAE=0.067，ATM相对误差0.31\%）和Heston（MAE=0.031，ATM相对误差0.52\%）上验证了统一框架的有效性，并与独立PINN基线进行了对比。
\end{itemize}

\section{论文结构}\label{sec:structure}

本文其余部分组织如下：第二章回顾期权定价模型、PINN方法和LLM辅助科学计算的相关工作；第三章介绍本文提出的统一PINN框架和LLM路由层的设计；第四章报告实验结果；第五章总结全文并讨论局限性与未来工作。
